\documentclass[10pt,a4paper]{amsart}
\usepackage{amsmath,amssymb,amsthm,graphicx}
% Retained minimal preamble for the Stacks Project chapter more-morphisms.tex (More on Morphisms).
% Verbatim excerpt of the chapter preamble: Xy-pic setup (native preamble lines 14-18) and the native
% theorem declarations with the mathematical macros (native preamble lines 154-217). Copied mechanically;
% nothing added, removed or redefined. The wrapper class and the amsmath/amssymb/amsthm/graphicx base
% packages are supplied by the assembler, so the class/externaldocument part of the native preamble is omitted.
% For commutative diagrams we use Xy-pic
\usepackage[all]{xy}

% We use 2cell for 2-commutative diagrams.
\xyoption{2cell}
% Theorem environments.
%
\theoremstyle{plain}
\newtheorem{theorem}[subsection]{Theorem}
\newtheorem{proposition}[subsection]{Proposition}
\newtheorem{lemma}[subsection]{Lemma}

\theoremstyle{definition}
\newtheorem{definition}[subsection]{Definition}
\newtheorem{example}[subsection]{Example}
\newtheorem{exercise}[subsection]{Exercise}
\newtheorem{situation}[subsection]{Situation}

\theoremstyle{remark}
\newtheorem{remark}[subsection]{Remark}
\newtheorem{remarks}[subsection]{Remarks}

\numberwithin{equation}{subsection}

% Macros
%
\def\lim{\mathop{\mathrm{lim}}\nolimits}
\def\colim{\mathop{\mathrm{colim}}\nolimits}
\def\Spec{\mathop{\mathrm{Spec}}}
\def\Hom{\mathop{\mathrm{Hom}}\nolimits}
\def\Ext{\mathop{\mathrm{Ext}}\nolimits}
\def\SheafHom{\mathop{\mathcal{H}\!\mathit{om}}\nolimits}
\def\SheafExt{\mathop{\mathcal{E}\!\mathit{xt}}\nolimits}
\def\Sch{\mathit{Sch}}
\def\Mor{\mathop{\mathrm{Mor}}\nolimits}
\def\Ob{\mathop{\mathrm{Ob}}\nolimits}
\def\Sh{\mathop{\mathit{Sh}}\nolimits}
\def\NL{\mathop{N\!L}\nolimits}
\def\CH{\mathop{\mathrm{CH}}\nolimits}
\def\proetale{{pro\text{-}\acute{e}tale}}
\def\etale{{\acute{e}tale}}
\def\QCoh{\mathit{QCoh}}
\def\Ker{\mathop{\mathrm{Ker}}}
\def\Im{\mathop{\mathrm{Im}}}
\def\Coker{\mathop{\mathrm{Coker}}}
\def\Coim{\mathop{\mathrm{Coim}}}

% Boxtimes
%
\DeclareMathSymbol{\boxtimes}{\mathbin}{AMSa}{"02}

%
% Macros for moduli stacks/spaces
%
\def\QCohstack{\mathcal{QC}\!\mathit{oh}}
\def\Cohstack{\mathcal{C}\!\mathit{oh}}
\def\Spacesstack{\mathcal{S}\!\mathit{paces}}
\def\Quotfunctor{\mathrm{Quot}}
\def\Hilbfunctor{\mathrm{Hilb}}
\def\Curvesstack{\mathcal{C}\!\mathit{urves}}
\def\Polarizedstack{\mathcal{P}\!\mathit{olarized}}
\def\Complexesstack{\mathcal{C}\!\mathit{omplexes}}
% \Pic is the operator that assigns to X its picard group, usage \Pic(X)
% \Picardstack_{X/B} denotes the Picard stack of X over B
% \Picardfunctor_{X/B} denotes the Picard functor of X over B
\def\Pic{\mathop{\mathrm{Pic}}\nolimits}
\def\Picardstack{\mathcal{P}\!\mathit{ic}}
\def\Picardfunctor{\mathrm{Pic}}
\def\Deformationcategory{\mathcal{D}\!\mathit{ef}}
\usepackage{cleveref}
\crefname{originallemma}{original lemma}{original lemmas}
\Crefname{originallemma}{Original lemma}{Original lemmas}
\makeatletter
\newlabel{lemma-control-agreement}{{lemma-control-agreement}{source}{}{sourcepointer-1}{}}
\newlabel{lemma-control-agreement@cref}{{[originallemma][0][]lemma-control-agreement}{[1][0][]source}{}{}{}}
\makeatother
\crefname{originallemma}{original lemma}{original lemmas}
\Crefname{originallemma}{Original lemma}{Original lemmas}
\makeatletter
\newlabel{more-algebra-lemma-approximate-complex-graded}{{07VF}{source}{}{sourcepointer-2}{}}
\newlabel{more-algebra-lemma-approximate-complex-graded@cref}{{[originallemma][0][]07VF}{[1][0][]source}{}{}{}}
\makeatother
\crefname{originallemma}{original lemma}{original lemmas}
\Crefname{originallemma}{Original lemma}{Original lemmas}
\makeatletter
\newlabel{more-algebra-section-artin-rees}{{07VD}{source}{}{sourcepointer-3}{}}
\newlabel{more-algebra-section-artin-rees@cref}{{[originallemma][0][]07VD}{[1][0][]source}{}{}{}}
\makeatother
\title{Conormal algebras of the fibres induced by a formal local map are controlled by one finite type modification}
\author{The Stacks Project Authors}
\date{stacks-project:more-morphisms.tex:local-snapshot-20261011; GNU FDL 1.2 or later}
\begin{document}
\maketitle
\noindent\textit{Editorial note.} Statement span is the whole lemma environment lines 11570-11590 including the displayed pair of closed immersions; proof span is lines 11592-11718.\par\medskip
\noindent\textit{Editorial note.} The proof cites several labelled results, which are prerequisites and are supplied as non-load-bearing navigation stubs.\par\medskip
\noindent\textit{Editorial note.} The author's prose contains an informal interjection; the text is retained verbatim.\par\medskip
\noindent\textit{Editorial note.} extra\_packages is empty on purpose: the retained preamble already loads xy with options and a second declaration clashes.\par\medskip
\noindent\textit{Allowed standard prerequisites (editorial).}\begin{itemize}
\item Explicit prerequisite boundary. Noetherian schemes of finite type, completions of local rings at their maximal ideal, base change of a morphism along a formal local map, and conormal algebras of closed immersions cut out by such base changes. The control construction and the comparison of conormal algebras are retained.
\end{itemize}\medskip
\ifdefined\hypertarget\hypertarget{sourcepointer-1}{}\fi
\noindent\textit{Source reference.} Original lemma lemma-control-agreement: sources/stacks-more-morphisms/pinned-original-source.tex, source lines 11530--11542. Referenced by the retained statement or proof but not part of the collected spans; supplied only as a navigation pointer. The argument does not depend on its content.\par\medskip
\ifdefined\hypertarget\hypertarget{sourcepointer-2}{}\fi
\noindent\textit{Source reference.} Original lemma 07VF: sources/stacks-more-morphisms/pinned-original-source.tex, source lines 11718--11718. Referenced by the retained statement or proof but not part of the collected spans; supplied only as a navigation pointer. The argument does not depend on its content.\par\medskip
\ifdefined\hypertarget\hypertarget{sourcepointer-3}{}\fi
\noindent\textit{Source reference.} Original lemma 07VD: sources/stacks-more-morphisms/pinned-original-source.tex, source lines 11718--11718. Referenced by the retained statement or proof but not part of the collected spans; supplied only as a navigation pointer. The argument does not depend on its content.\par\medskip
\begin{lemma}
\label{lemma-control-graded}
Let $Y \to T \to S$ be finite type morphisms of Noetherian schemes.
Let $t \in T$ map to $s \in S$ and let
$\sigma : \mathcal{O}_{T, t} \to \mathcal{O}_{S, s}^\wedge$
be a local $\mathcal{O}_{S, s}$-algebra map.
There exists a finite type morphism $(T', t') \to (T, t)$
such that $\sigma$ factors through
$\mathcal{O}_{T, t} \to \mathcal{O}_{T', t'}$
and such that for every local $\mathcal{O}_{S, s}$-algebra map
$\sigma' : \mathcal{O}_{T, t} \to \mathcal{O}_{S, s}^\wedge$
which factors through $\mathcal{O}_{T, t} \to \mathcal{O}_{T', t'}$
the closed immersions
$$
Y \times_{T, \sigma} \Spec(\mathcal{O}_{S, s}^\wedge) = Y_\sigma
\longleftarrow Y_t \longrightarrow
Y_{\sigma'} =
Y \times_{T, \sigma'} \Spec(\mathcal{O}_{S, s}^\wedge)
$$
have isomorphic conormal algebras.
\end{lemma}

\begin{proof}
A useful observation is that $\kappa(s) = \kappa(t)$ by the existence of
$\sigma$. Observe that the statement makes sense as the fibres of $Y_\sigma$
and $Y_{\sigma'}$ over $s \in \Spec(\mathcal{O}_{S, s}^\wedge)$
are both canonically isomorphic to $Y_t$. We will think of the property
``$\sigma'$ factors through $\mathcal{O}_{T, t} \to \mathcal{O}_{T', t'}$''
as a constraint on $\sigma'$. If we have several such constraints,
say given by $(T'_i, t'_i) \to (T, t)$, $i = 1, \ldots, n$
then we can combined them by considering
$(T'_1 \times_T \ldots \times_T T'_n, (t'_1, \ldots, t'_n)) \to
(T, t)$. We will use this without further mention in the following.

\medskip\noindent
By Lemma \ref{lemma-control-agreement} we can assume that any $\sigma'$
as in the statement of the lemma is the same as $\sigma$ modulo
$\mathfrak m_s^2$. Note that the conormal algebra of $Y_t$ in
$Y_\sigma$ is just the quasi-coherent graded $\mathcal{O}_{Y_t}$-algebra
$$
\bigoplus\nolimits_{n \geq 0}
\mathfrak m_s^n\mathcal{O}_{Y_\sigma}/
\mathfrak m_s^{n + 1}\mathcal{O}_{Y_\sigma}
$$
and similarly for $Y_{\sigma'}$. Since $\sigma$ and $\sigma'$
agree modulo $\mathfrak m_s^2$ we see that these two algebras
are the same in degrees $0$ and $1$. On the other hand, these
conormal algebras are generated in degree $1$ over degree $0$.
Hence if there is an isomorphism extending the isomorphism
just constructed in degrees $0$ and $1$, then it is unique.

\medskip\noindent
We may assume $S$ and $T$ are affine. Let $Y = Y_1 \cup \ldots \cup Y_n$
be an affine open covering. If we can construct $(T_i', t'_i) \to (T, t)$
as in the lemma such that the desired isomorphism (see previous paragraph)
exists for $Y_i \to T \to S$ and $\sigma$, then these glue by
uniqueness to prove the result for $Y \to T$. Thus we may assume $Y$ is affine.

\medskip\noindent
Write $S = \Spec(R)$, $T = \Spec(C)$, and $Y = \Spec(B)$.
Choose a presentation $B = C[x_1, \ldots, x_n]/(f_1, \ldots, f_m)$.
Denote $R^\wedge = \mathcal{O}_{S, s}^\wedge$.
Let $a_{kj} \in R^\wedge[x_1, \ldots, x_n]$ be polynomials
such that
$$
\sum\nolimits_{j = 1, \ldots, m} a_{kj}\sigma(f_j) = 0,\quad
\text{for }k = 1, \ldots, K
$$
is a set of generators for the module of relations among
the $\sigma(f_j) \in R^\wedge[x_1, \ldots, x_n]$.
Thus we have an exact sequence
\begin{equation}
\label{equation-resolution}
R^\wedge[x_1, \ldots, x_n]^{\oplus K} \to
R^\wedge[x_1, \ldots, x_n]^{\oplus m} \to
R^\wedge[x_1, \ldots, x_n] \to B \otimes_{C, \sigma} R^\wedge \to 0
\end{equation}
Let $c$ be an integer which works in the Artin-Rees lemma for
both the first and the second map in this sequence and the ideal
$\mathfrak m_{R^\wedge}R^\wedge[x_1, \ldots, x_n]$ as defined in
More on Algebra, Section \ref{more-algebra-section-artin-rees}.
Write
$$
a_{kj} = \sum\nolimits_{I \in \Omega} a_{kj, I} x^I
\quad\text{and}\quad
f_j = \sum\nolimits_{I \in \Omega} f_{j, I} x^I
$$
in multiindex notation where $a_{kj, I} \in R^\wedge$, $f_{j, I} \in C$,
and $\Omega$ a finite set of multiindices. Then we see that
$$
\sum\nolimits_{j = 1, \ldots, m,\ I, I' \in \Omega,\ I + I' = I''}
a_{kj, I} \sigma(f_{j, I'}) = 0,\quad
I''\text{ a multiindex}
$$
in $R^\wedge$. Thus we take
$$
C' = C[t_{jk, I}]/
\left(
\sum\nolimits_{j = 1, \ldots, m,\ I, I' \in \Omega,\ I + I' = I''}
t_{kj, I} f_{j, I'},\ I''\text{ a multiindex}\right)
$$
Then $\sigma$ factors through a map $\tilde\sigma : C' \to R^\wedge$
sending $t_{kj, I}$ to $a_{jk, I}$.
Thus $T' = \Spec(C')$ comes with a point $t' \in T'$ such that
$\sigma$ factors through $\mathcal{O}_{T, t} \to \mathcal{O}_{T', t'}$.
Let $t_{kj} = \sum t_{kj, I} x^I$ in $C'[x_1, \ldots, x_n]$.
Then we see that we have a complex
\begin{equation}
\label{equation-resolution-new}
C'[x_1, \ldots, x_n]^{\oplus K} \to
C'[x_1, \ldots, x_n]^{\oplus m} \to
C'[x_1, \ldots, x_n] \to B \otimes_C C' \to 0
\end{equation}
which is exact at $C'[x_1, \ldots, x_n]$ and whose base change
by $\tilde\sigma$ gives (\ref{equation-resolution}).

\medskip\noindent
By Lemma \ref{lemma-control-agreement}
we can find a further morphism $(T'', t'') \to (T', t')$
such that $\tilde\sigma$
factors through $\mathcal{O}_{T', t'} \to \mathcal{O}_{T'', t''}$
and such that if $\sigma' : C \to R^\wedge$
factors through $\mathcal{O}_{T'', t''}$, then the induced map
$\tilde \sigma' : C' \to R^\wedge$
agrees modulo $\mathfrak m_s^{c + 1}$ with $\tilde \sigma$.
Thus if $\sigma'$ is such a map, then we obtain a complex
$$
R^\wedge[x_1, \ldots, x_n]^{\oplus K} \to
R^\wedge[x_1, \ldots, x_n]^{\oplus m} \to
R^\wedge[x_1, \ldots, x_n] \to B \otimes_{C, \sigma'} R^\wedge \to 0
$$
over $R^\wedge[x_1, \ldots, x_n]$ by applying $\tilde\sigma'$
to the polynomials $t_{kj}$ and $f_j$. In other words, this
is the base change of the complex (\ref{equation-resolution-new})
by $\tilde\sigma'$. The matrices defining this complex
are congruent modulo $\mathfrak m_s^{c + 1}$ to the matrices
defining the complex (\ref{equation-resolution}) because
$\tilde \sigma$ and $\tilde \sigma'$ are congruent modulo
$\mathfrak m_s^{c + 1}$. Since (\ref{equation-resolution}) is exact,
we can apply
More on Algebra, Lemma \ref{more-algebra-lemma-approximate-complex-graded}
to conclude that
$$
\text{Gr}_{\mathfrak m_s}(B \otimes_{C, \sigma'} R^\wedge)
\cong
\text{Gr}_{\mathfrak m_s}(B \otimes_{C, \sigma} R^\wedge)
$$
as desired.
\end{proof}

\end{document}
